\chapter{Dysphagia Tests}
\section*{What Is This Test?} Dysphagia tests evaluate difficulty moving food or liquid from the mouth to the stomach and distinguish oropharyngeal from esophageal causes.
\section*{Alternative Names} Swallowing evaluation; swallowing study; esophageal evaluation.
\section*{Principle} Clinical assessment, modified barium swallow, barium esophagram, endoscopy, and manometry examine different phases and mechanisms.
\section*{Why This Test Is Ordered} It is used for choking, coughing during meals, food sticking, weight loss, aspiration, pain with swallowing, or recurrent pneumonia.
\section*{Primary vs Secondary Test} History and bedside swallowing assessment are primary; imaging, endoscopy, and physiologic studies are selected secondary tests.
\section*{How This Test Is Performed} Depending on the question, the patient may swallow food or contrast under fluoroscopy, undergo endoscopy, or have pressure measured with manometry.
\section*{What Preparation Is Needed} Follow fasting instructions for imaging or endoscopy. Bring a medication and symptom history.
\section*{Understanding Results} Findings may show aspiration, weakness, obstruction, narrowing, reflux, or abnormal motility. A normal test does not exclude intermittent symptoms.
\section*{Follow-up Tests} Speech therapy, endoscopy with biopsy, manometry, pH monitoring, CT, or neurologic evaluation may follow.
\section*{References} \begin{enumerate}\item MedlinePlus. Dysphagia Tests.\item American Speech-Language-Hearing Association. Adult dysphagia assessment guidance.\end{enumerate}
\clearpage\section*{Understanding Results and Follow-up}
The laboratory report should identify the specimen, method, units, reference interval or decision limit, and whether the result is positive, negative, abnormal, or inconclusive. Reference intervals vary with age, sex, pregnancy, specimen type, assay, and laboratory; a result outside a range is not automatically a diagnosis, and a result within a range does not always exclude disease. Discuss unexpected results with the ordering clinician, who can decide whether repeat or confirmatory testing, treatment, or referral is appropriate. Do not change medicines or delay urgent care based on an isolated result.
\section*{Regulatory Status and Authorization Date}This chapter describes a test category, not a specific commercial product. FDA, CDSCO, EU IVDR, and MHRA approval or registration dates therefore cannot be assigned without the assay manufacturer, product name, instrument, and intended use. For laboratory-developed tests, an individual agency approval date may not exist. See the Preface for the interpretation of regulatory status.
\clearpage
